\section{Model and assumptions}
\label{sec:model}
\subsection{L1 control state}
\label{sec:l1}
Nodes authenticate finalized L1 headers, authorities and a complete
ordered control prefix. L1 timestamps label L1 blocks; they are not
regional commit times or Espresso activation times.

\begin{table}[htbp]\centering\small
\caption{L1 authority and control records.}\label{tab:l1}
\begin{tabularx}{\linewidth}{@{}p{.27\linewidth}Y@{}}\toprule
Field & Meaning\\\midrule
$\tagg{AUTH}_{r,0}$ & The fixed regional snapshots, leader schedule,
DATA quota $\TB_r$, and export-window parameters for index $r$.\\
$\tagg{AUTH}_{r,g}$, $g>0$ & The leader schedule for recovery
generation $g$; it does not change the original snapshot or quota.\\
Control sequence $c$ & A unique increasing index in the L1 control
log. It binds the chain, record kind and payload.\\
Activation time $\theta_a$ & An Espresso timestamp threshold,
announced before it can be crossed.\\
STATUS & A future live/off update for the named regions. It changes
selection eligibility, not existing voting weights.\\
MSIG FORCE & A target branch version, activation time and any authorized
status changes to apply with the Force. No retained height is supplied.\\
CLEAR & An explicit future removal of the named automatic cooldown;
ordinary ON does not perform this operation.\\\bottomrule
\end{tabularx}
\end{table}

\paragraph{Authority records.}
For each $(r,g)$, L1 supplies at most one immutable authority.
$\tagg{AUTH}_{r,0}$ fixes $\mathit{cfg}_r$, including the members,
weights, $F_i$, $\Delta_i$ and $K$ of each region, and fixes $\TB_r$.
It also fixes positive window lengths $W_{\rm start},W_{\rm run}$,
positive count targets $m_{\rm start},q_{\min}$, and cooldown
$C_{\rm auto}>0$. Counts are capped by the remaining DATA quota.
For $g>0$, only the leader schedule changes.
Nodes wait for required authority evidence. Later controls and slashing
never reweight an old certificate or reset a session quota.

Authority is indexed by $r$, not by a particular speculative branch.
A new branch version that reuses $r$ uses the same authority at that index.
Its pair $(\beta,r)$ identifies a different engine. Rebinding an already
assigned authority is not permitted.

\paragraph{Control ingestion and activation.}
The complete L1 control prefix has a sequence number and a binding
history commitment. A verifier retains pending records as well as
already applied records. Ingestion does not imply activation.
Activation times need not follow sequence order: records becoming due
together are ordered by $(\theta_a,c)$. Malformed, unauthorized or
duplicate records do not grant a second action.

STATUS and CLEAR are independent controls. Configuration attached to
a Force applies only if that Force's target branch version is still current.
A stale Force is consumed without changing the new branch version. An independent
STATUS beside it still takes effect on its own terms.
Records are immutable after finalization. A revision is a new,
prospectively scheduled record, not an edit to past history.

\paragraph{Espresso positions and effective status.}
Let $J_q$ be Espresso block $q$, with authenticated timestamp
$\theta_q$. Positions also include namespace-transaction and record
indices. Heights and those indices give a total order even when
timestamps are equal. For a threshold $d$, define
\[
 q(d)=\min\{q:\theta_q\ge d\}.
\]
Its event is applied at the opening of $J_{q(d)}$, before that block's
EXPORTs. The prelude applies due independent status changes, then a
matching multisignature Force, then any automatic Force.

Write $\chi_E(q)$ for the eligibility snapshot after this prelude and
before the EXPORTs. It combines the L1 live/off map with automatic
cooldowns through that position. A region is eligible only if its
governance flag is on, its cooldown has expired, and its designated
authority has positive total weight. The historical snapshot is
authenticated by replay, not by an endpoint's statement.

OFF changes subsequent selections made from the affected snapshots.
It does not revoke an existing historical REGULAR witness. ON does not
clear cooldown or reopen an old branch version. A committee that needs to stop
the current branch uses a Force. No control changes an accepted DATA
prefix or retroactively changes a regional voting snapshot.

\paragraph{Genesis.}
The initial checkpoint fixes the chain identifier, cyclic order
$\mathcal L_R$, genesis DATA anchor, first replay position, Espresso
quorum trust anchor, control prefix and initial branch version $\beta=0$.
A synthetic boundary selects the first eligible region in
$\mathcal L_R$ at the first eligible authenticated position. It waits
if none exists. The initial authority and state are trusted or verified
from their specified sources; Espresso does not authenticate earlier
application history merely by starting this interpreter.
The genesis anchor includes synthetic commitment evidence and fixed,
authenticated L1 and Espresso references. Its pending boundary has
source context $(i_b,\beta_b,r_b)=(\bot,\bot,-1)$. It opens
session $(\beta,0)$ under the version current at selection: normally
$(0,0)$, or a later version if a Force acted while genesis was Waiting.
At selection height $q_s$, set $q^{\rm open}_{\beta,0}=q_s$ and
$\lambda^{\rm open}_{\beta,0}=\max\{\lambda(\text{genesis}),
\lambda(\tagg{AUTH}_{0,0})\}$, where the second term is the authority's
L1 publication height. A supplied later checkpoint must instead include
its current boundary, reference bounds and complete interpreter state.

\subsection{Data delivery and the Publish Layer}
\label{sec:delivery}
The \defterm{Publish Layer} is a delivery service implemented by existing
participants. It does not run another consensus protocol or choose which
DATA is valid. A regional leader, another replica, or an external verifier
can perform a submission. The rotation rules use the service guarantee
below; \cref{sec:publish-implementation} describes one implementation.

A \defterm{complete regional result} is DATA with its direct or covering
FINAL receipt and the ancestry, boundaries and source evidence needed to
validate it. Correct replicas offer complete results to the Publish Layer
as soon as they obtain them. The service may batch results and construct
EXPORT records as specified in \cref{sec:exports}. A proposal QC, a block
hash, or a receipt whose payload is missing is not a complete input.

Regions may send the same results directly to other regions. This lets a
regular successor reach Ready before the terminal result is finalized in
Espresso. If direct delivery is unavailable, it can obtain the result from
Espresso after the record is finalized. The latter path waits for Espresso; the
protocol does not impose that wait when direct delivery succeeds.
There is no separate cross-region transport service in this model.

Espresso carries the finalized records, including DATA and commitment
evidence, to online readers. Readers still check the application and
rotation rules. L1 supplies its own authority and control proofs.
An execution-state snapshot may be transferred directly or reconstructed
from authenticated input history. State reconstruction is part of Ready,
not something that Espresso inclusion performs automatically.

\subsection{Stake and faults}
\label{sec:stake}
Member $v$ of region $i$ has voting weight $\omega_{i,v}>0$. Let
$W_i=\sum_{v\in V_i}\omega_{i,v}$. The stake configuration
fixes a bound $F_i<W_i/2$, and Byzantine members of region $i$ hold
total weight at most $F_i$. Every regional certificate needs
matching signatures from distinct members with total weight greater than
$F_i$. So every certificate has a correct signer, and the correct members
alone can form one. This replaces the $f+1$ of $n=2f+1$ in Part~I.

Membership and weights stay fixed for a whole session, including its
recoveries. Slashing and status updates do not change them. A later
session checks old certificates against the old snapshot. The bound
$F_i$ counts every key of the snapshot that is compromised while its
certificates are still accepted. Compromise of old keys after that is
outside this model.

Correct replicas keep their protocol state and committed blocks.
Signatures cannot be forged, and hashes are binding.



\subsection{Assumptions}
\label{sec:assumptions}
\begin{assumption}[Finalized L1 records]
\label{ass:authority}
At every finalized L1 height, the authority records and ordered control
history through that height are common and authenticated. Finalized
records never change. A verifier obtains a complete control prefix,
not an arbitrary subset of valid events.
\end{assumption}

\begin{assumption}[Authenticated synchronous communication and global clock]
\label{ass:network}
\leavevmode\par
\begin{itemize}
\item Messages use authenticated communication channels.
\item Within region $i$, the actual delay of a correct-to-correct
message is $\delta\le\Delta_i$; $\delta$ may differ across messages.
\item All nodes share a global clock.
\end{itemize}
\end{assumption}
The delay bound includes delivery and processing of complete validation
dependencies. Each replica records its own entry and voting events on
the global clock.

\begin{assumption}[Bounded publication service]
\label{ass:publish}
If a correct participant has a complete regional result available to the
Publish Layer at real time $s$, the service submits an Espresso-admissible
record carrying that result by $s+T_{\rm pub}$. The service preserves its
DATA and commitment evidence and does not depend on the original leader
alone. At least one working submission path and the required fee resources
are available. Here $T_{\rm pub}>0$ is fixed for the stated input-size and load
limits. A result already recorded in Espresso satisfies this delivery
guarantee; retries needed for acceptance into $G$ are a separate task.
\end{assumption}
The clock starts when the complete result is available, not when a leader
first proposes a block or one replica first commits it. Forming the FINAL
receipt remains a regional-protocol task. ``Submits'' means delivery of the
complete record to the Espresso submission interface under its admission
rules, not merely a local send attempt.


\begin{assumption}[Honest multisignature committee]
\label{ass:committee}
The committee is one honest, coordinated control service. It submits
only authorized, immutable, future-dated STATUS, CLEAR and MSIG FORCE
records. It chooses activation times satisfying the advance-announcement
condition below. It keeps no private protected prefix and has no power
to select an older retained anchor.
\end{assumption}

\begin{assumption}[Bounded L1 finalization and visibility]
\label{ass:monitor}
A valid control operation actually submitted at real time $s$ is
finalized, available and verifiable by every correct online participant
by $s+T_{\rm L1}$. Participants include regional replicas, external
observers and certifiers. The bound includes inclusion, finalization,
delivery, retrieval and validation.
\end{assumption}

\Needspace{8\baselineskip}
\begin{assumption}[Espresso history and time]
\label{ass:espresso}
Correct verifiers obtain compatible finalized Espresso prefixes and
can authenticate each required header and complete namespace payload,
or its genuine absence. The history and evidence needed for replay remain
retrievable. In particular, globally accepted DATA remains recoverable
without the original producing region staying online. Timestamps are
non-decreasing. If block $J_q$ first becomes final at real time $t_E(q)$,
then $\theta_q\le t_E(q)+\epsilon_E$, where $\epsilon_E\ge0$ is the
configured maximum clock lead.
\end{assumption}
Long-term retrieval is an explicit service assumption. It may be supplied
by archival readers that preserve authenticated Espresso history; it does
not follow from a finality certificate alone. A returning node fetches
missing history or starts from an authorized checkpoint before confirming
new output. Catching up with arbitrarily long history is not assigned a
constant-time bound.

\begin{assumption}[Bounded Espresso finalization and visibility]
\label{ass:espresso-visibility}
An admissible complete record submitted to Espresso at real time $s_E$
is included in a finalized block by $s_E+T_E$, for a fixed $T_E>0$.
By the same deadline, every
correct participant that stays online and keeps up with the history can
retrieve the complete record and verify its inclusion and finality.
The bound covers inclusion, finalization, dissemination, retrieval and
verification of the Espresso evidence, under the declared admission and
load conditions.

Each newly finalized Espresso block and its complete relevant namespace
also become available and verifiable to those participants within $T_E$
of finalization, including blocks with no records for this chain.
\end{assumption}
The service provides full payloads and evidence, not only header hashes.
It does not promise that an EXPORT passes this chain's acceptance rules:
a stale base, closed branch or missing predecessor can still prevent it
from advancing $G$. Regional and application validation remain separate.

\paragraph{The two delivery bounds.}
For a newly available complete result, the two assumptions give
\[
 s_{\rm submit}\le s+T_{\rm pub},\qquad
 s_{\rm visible}\le s+T_{\rm pub}+T_E.
\]
Here $s_{\rm visible}$ means that the submitted result and its Espresso
evidence are available to the online readers above. It is not a bound on
when the regional result first forms, when a new region restores its
execution state, or when a rejected EXPORT is accepted after retries.
These are real-time bounds. Delivery windows use Espresso timestamps, so
window parameters also need the clock profile and regional progress and
startup allowances. The bound $T_E$ does not replace the clock-lead
assumption or the advance-announcement rule below.


\paragraph{Advance announcement.}
For a control submitted at $s$, the committee chooses
\[
 \theta_a>s+T_{\rm L1}+\epsilon_E+\eta,
 \qquad \eta>0.
\]
Thus the control is available to correct participants before its
first crossing can become final. This is a model condition, not a
conclusion from an average L1 or Espresso block time. The L1 service
bound already includes dissemination; no second announcement delay is
added. Offline or unsynchronized participants catch up before confirming
new output. Missing control history pauses verification.

\paragraph{Verification interfaces.}
Regional receipt validity and the application-validity checks below
are assumed sound under the stated regional model. A deployed
certifier must authenticate both source histories and bind the control
version used. A signature or TEE alone does not establish that an old
control snapshot is complete. The publication interface additionally
requires sound execution and DA correspondence checks.

\subsection{Conditions for progress}
\label{sec:progress}
The following conditions are separate from prefix safety.
\begin{itemize}
\item Correct leaders receive repeated turns that are actually entered,
including recovery initializers and DATA-producing turns across
sessions. Exposed equivocators are eventually excluded permanently
from later leader schedules.
\item Espresso continues to finalize retrievable history with growing
timestamps and includes this chain's valid exports.
\item Complete results become available to the Publish Layer. Its
submission paths and Espresso meet their stated bounds. Retried EXPORTs
for the next required DATA are accepted in time during some sustained
service interval.
\item Required authorities and recoverable native state become
available. A waiting boundary eventually has another eligible region.
The committee leaves working regions enough uninterrupted service time.
\item Correct clients' admissible requests reach the regions that
continue serving. Requests fit in DATA batches and have only finitely
many older requests ahead of them in each local queue.
\end{itemize}
The bounds begin with a complete result and do not bound regional
recovery, state restoration, or repeated rejected EXPORTs. Window success
therefore still needs the stated serving interval and a compatible
Espresso clock profile. Otherwise repeated automatic Forces can prevent
progress without revoking any global DATA. Espresso failure also stops external time and
global progress. Switching the global finality service is outside this
protocol.

